% Generated draft: requires mathematical and source review.
\documentclass[11pt]{article}
\usepackage[margin=25mm]{geometry}
\usepackage{fontspec}
\setmainfont{Latin Modern Roman}
\usepackage{amsmath,amssymb,mathtools,unicode-math}
\setmathfont{Latin Modern Math}
\usepackage{enumitem,needspace,xurl,hyperref}
\setlength{\parindent}{0pt}
\setlength{\parskip}{5pt}
\setlength{\emergencystretch}{3em}
\newcommand{\N}{\mathbb{N}}
\newcommand{\Z}{\mathbb{Z}}
\newcommand{\Q}{\mathbb{Q}}
\newcommand{\R}{\mathbb{R}}
\newcommand{\C}{\mathbb{C}}
\newcommand{\F}{\mathbb{F}}
\newcommand{\A}{\mathbb{A}}
\newcommand{\PP}{\mathbb{P}}
\newcommand{\E}{\mathbb{E}}
\newcommand{\eps}{\varepsilon}
\newcommand{\dd}{\,\mathrm{d}}
\begin{document}
\section*{UnsolvedMath ID: 30005296}
\textbf{Source problem code:} catalogue entry
\subsection{Definitions and mathematical statement}
The supplied record does not specify a ground field or its characteristic. Write \(k\) for a ground field whose hypotheses remain to be determined. Here, ``general'' means belonging to a nonempty Zariski open subset of the relevant parameter space; whether this is to be interpreted geometrically or using \(k\)-rational parameters also requires clarification.

For positive integers \(a,b\), a finite reduced point set \(Z\subset\PP^3_k\) is called \((a,b)\)-geproci if its projection from a general admissible center \(p\in\PP^3_k\setminus Z\) to a projective plane is injective on \(Z\), and the projected reduced scheme is a scheme-theoretic complete intersection of plane curves of degrees \(a\) and \(b\). Thus, for such a projection \(\pi_p\), there are homogeneous forms \(F_p,G_p\) of degrees \(a,b\) with no common curve component such that
\[
\pi_p(Z)=V(F_p,G_p)\subset\PP^2_k,
\qquad |Z|=ab.
\]
Call \(Z\) geproci if it is \((a,b)\)-geproci for some positive integers \(a,b\). These definitions state the usual finite reduced point-set interpretation; the intended field, rationality conventions, and any additional restrictions must be checked against the source.

An \((r,s)\)-grid in \(\PP^3_k\) consists of the \(rs\) distinct intersection points of two families of distinct lines
\[
L_1,\ldots,L_r,\qquad M_1,\ldots,M_s,
\]
where the lines within each family are pairwise skew and every \(L_i\) meets every \(M_j\). A finite point set is in linear general position in \(\PP^3_k\) if every subset of \(t\leq4\) points spans a projective subspace of dimension \(t-1\).

The supplied original excerpt presents the following related problems.
\begin{enumerate}
\item Classify geproci point sets in \(\PP^3\), with the \((4,4)\)-geproci case identified as a starting point. The cleaned statement restricts the classification to ``nontrivial'' sets, but does not define that restriction or the equivalence relation for classification.
\item Determine whether every geproci set \(Z\subset\PP^3\) with \(|Z|\geq5\) contains three distinct collinear points; equivalently, whether there exist distinct \(x,y,z\in Z\) and a line \(\ell\subset\PP^3\) such that \(x,y,z\in\ell\). The excerpt states this question for every geproci set, whereas the surrounding discussion refers to nontrivial sets. The intended restriction is unresolved. The excerpt describes an affirmative answer as implying that no nontrivial geproci point set is in linear general position; this is source-supplied motivation, with its intended size range and meaning of nontriviality requiring clarification.
\item As a proposed first step, determine whether, for positive integers \(a,d\), a smooth curve \(C\subset\PP^3\) of degree \(d\), and a set \(Z\subset C\) of \(ad\) distinct points in linear general position, \(Z\) can ever be \((a,d)\)-geproci. The original excerpt proposes the universal negative statement
\[
\bigl(C\subset\PP^3\text{ smooth},\ \deg C=d,\ Z\subset C,\ |Z|=ad,\ Z\text{ in linear general position}\bigr)
\Longrightarrow Z\text{ is not }(a,d)\text{-geproci}.
\]
This is a proposed target, not an asserted theorem. Missing restrictions on \(a,d,C\), and on nontriviality, remain unresolved. The cleaned statement instead asks about a general finite subset of a smooth plane curve. A smooth space curve with a linear-general-position subset and a smooth plane curve with a general subset are different formulations; the supplied evidence does not justify identifying them. In particular, neither the cardinality nor the parameter space meant by ``general finite subset'' is specified in the cleaned version.
\item Determine whether every \((a,b)\)-geproci set contains an \((a',b')\)-grid satisfying the quantitative lower bounds intended in the source. The excerpt says that every known example contains such a grid, and presents this proposed property as considerably stronger than the linear-general-position question. Its extracted bounds read ``\(a'\geq\) \texttt{a2}'' and ``\(b'\geq\) \texttt{2b}''. The typography is not sufficiently reliable to recover the intended bounds; no corrected inequalities are asserted here. Consequently, the quantitative grid question remains incompletely specified.
\end{enumerate}
The excerpt also motivates characterizing subschemes of \(\PP^N\) having the geproci property. It does not supply the scheme-theoretic definition or projection conventions needed to formulate that broader target, so the point-set definition above is not asserted to cover it.
\subsection{Short English statement}
Classify the intended nontrivial geproci point sets in projective three-space, beginning with the \((4,4)\) case, and formulate the proposed collinearity, smooth-curve, and contained-grid questions. A geproci set has a general planar projection that is a complete intersection. The field, nontriviality restriction, curve formulation, and extracted grid bounds require clarification before all targets can be stated unambiguously.
\subsection{Sources}
\begin{itemize}
\item Supplied local problem-registry record: permanent UnsolvedMath ID \(30005296\), title \emph{Classification and Grid Structure of Geproci Sets}, exact source problem code \texttt{catalogue entry}. The project identifies \texttt{data/unsolvedmath/problems.json} as its primary registry and \texttt{catalog/problems.json} as a supplemental registry; this draft uses only the supplied complete record.
\item Supplied report citation: \emph{Mini-Workshop: Subvarieties in Projective Spaces and Their Projections (2023)}, Oberwolfach Reports, pp.\ 3093--3119, DOI \texttt{10.4171/owr/2022/54}. \url{https://doi.org/10.4171/owr/2022/54}. The registry gives proposed year 2022 and the tag \texttt{owr-year:2022}, while its report citation names 2023. These supplied dates are retained without resolving their relationship. No individual proposer is supplied.
\item Supplied literature reference: Chiantini et al., \emph{Classification of certain geproci sets}. \url{https://arxiv.org/abs/2303.16263}.
\item Supplied literature reference: Chiantini et al., \emph{Geproci sets and skew-line combinatorics}. \url{https://arxiv.org/abs/2308.00761}.
\item Supplied literature reference: EMS Press, \emph{OWR 2022/54}. \url{https://doi.org/10.4171/owr/2022/54}.
\item Literature assessment dated exactly \texttt{2026-08-22}: the supplied curation labels the problem \texttt{partially\_solved} and reports progress through constructions and classifications of several families, including the 16-point four-skew-line case, while reporting that no classification of all geproci sets is known. This is supplied dated curation, not independently verified here. The background explicitly requests independent verification of time-sensitive status and cited claims.
\item The registry labels the statement \texttt{corrected\_verified} and says that it was checked against the official Oberwolfach report. This is a supplied verification claim, not a verification performed for this draft; the discrepancies and corrupted notation described above remain unresolved.
\end{itemize}
\end{document}
